\section{The scaling of $\mu$}
\label{runtimescalingofUpsi}
The asymptotic behavior of runtime, discussed in the main text, is governed by the scaling of $\mu$ \eqref{mudefnition} with dimension $d$ of the Hilbert space. The dependence $\mu(d)$ is specific to the matrix $A$ to be exponentiated (which is proportional to the Hamiltonian). This prevents us from stating a general scaling relation. We thus turn to a discussion of several concrete examples. 

We first consider the example of the transmon-cavity system, see Eq.\ \eqref{fockH}, where we increase $d$ via the cavity dimension $d_c$ while leaving the transmon cutoff fixed. The corresponding Hamiltonian $H_1$, written in the bare product basis, then has a constant maximum number of non-zero elements in each row. (I.e., the parameter $\sigma'$ \eqref{sigmapdef} is independent of $d$.) Numerically, we find in this case that $\mu$ scales linearly with $\norm{A}_1$ \big($\mu\sim\Theta(\norm{A}_1)$\big), and this linearity holds for a wide range of $\tau$, see Fig.\ \ref{fig:muscaling}(a). In this example, two ingredients lead to the scaling $\norm{A}_1\sim \Theta(d^{\frac{1}{2}})$ : (1) $d$ is a constant multiple of $d_c$, and (2) the cavity ladder operators in $H_1$ which has one-norm scaling $\Theta(d_c^\frac{1}{2})$. Thus, we obtain the overall scaling $\mu\sim \Theta(d^{\frac{1}{2}})$, see Fig.\ \ref{fig:muscaling}(e). In the second example, we consider the system three coupled transmons, see Eq.\ \eqref{gateH}. The Hamiltonian $H_2$ is written in the harmonic oscillator basis and has a constant $\sigma'$, which leads to $\mu\sim\Theta(\norm{A}_1)$ [Fig.\ \ref{fig:muscaling}(b)]. One-norm of operator $(b^\dagger b)^2$ in $H_2$ has quadratic scaling with each transmon dimension. Since we increase $d$ by enlarging the dimensions of all three transmons simultaneously, we obtain $\norm{A}_1\sim\Theta(d^{\frac{2}{3}})$. This scaling results in $\mu\sim \Theta(d^{\frac{2}{3}})$, see Fig.\ \ref{fig:muscaling} (f).

For other Hamiltonian matrices where $\sigma'$ depends on $d$, the scaling $\mu\sim\Theta(\norm{A}_1)$ may still hold. In the following, we show two concrete examples that one where linearity holds, and another where it breaks. We first consider a driven system of $N_q$ qubits with nearest-neighbor $\sigma_z$ coupling. In the rotating frame, the system is described by 
\begin{align*}
 H_3(t)=\sum_{\nu=1}^{N_q}[a_x^{(\nu)}(t)\sigma^{(\nu)}_x+a_y^{(\nu)}(t)\sigma^{(\nu)}_y]+\sum_{\nu=1}^{N_q-1}g\sigma_z^{(\nu)}\sigma_z^{{(\nu+1)}}.
\end{align*}
\begin{figure}
  \includegraphics[width=1\linewidth]{uscaling.pdf}

  \caption{Scaling of $\mu$, the number of matrix-vector multiplications. (a)-(d) The scaling of $\mu$ with $\norm{A}_1$ for Hamiltonians $H_j$, $j=1,2,3,4$, when considering several error tolerances $\tau$. (e)-(h) The scaling of $\mu$ with dimension $d$ when considering $\tau=10^{-8}$ for each $H_j$, respectively.  In all figures, dots are sampling data and solid lines are fitting curves. The parameters setup of $H_1,\ H_2$ are the same as the ones given in the captions of Figs.\ \ref{fig:state_time}, \ref{fig:gate}. Parameters of $H_3$ and $H_4$ are as follows: for $H_3$, $E_{Ci}/2\pi=2.5$\ GHz, $E_{Ji}/2\pi=8.9$\ GHz, $E_{Li}/2\pi=0.5$\ GHz, $g/2\pi=0.1\ $GHz and $\phi=0.33$ ; for $H_4$, $a_x^{(\nu)}/2\pi=a_y^{(\nu)}/2\pi=0.5$\ GHz and $g/2\pi=0.1$\ GHz. }
  \label{fig:muscaling}
\end{figure}
Each qubit is controlled via drive fields $a_x^{(n)}$ and $a_y^{(n)}$. 
For this example, the numerical results show that scaling $\mu\sim\Theta(\norm{A}_1)$ still holds, see Fig.\ \ref{fig:muscaling}(c). Since we increase $d$ by enlarging $N_q$, we obtain $\norm{A}_1\sim\Theta(N_q)\sim\Theta(\log_2d)$. This scaling results in $\mu\sim \Theta(\log_2d)$, see Fig.\ \ref{fig:muscaling}(g). For the second example, we consider two capacitively coupled fluxonium qubits with Hamiltonian
\begin{align}
\label{fluxonium}    H_4(t)&=\sum_{i=1}^2 \big(4E_{Ci}n_i^2-E_{Ji}\cos(\varphi_i)+\frac{E_{Li}}{2}[\varphi_i-\phi_i(t)]^2\big) + gn_1n_2.
\end{align}
Here, $E_{Ci}$, $E_{Ji}$, $E_{Li}$ denote the charging, inductive and Josephson energies of two fluxonium qubits, and $g$ is the interaction strength. $\varphi$ and $n$ are the phase and charge number operators, defined in the usual way. We represent the Hamiltonian in the harmonic oscillator basis. Fig.\ \ref{fig:muscaling}(d) shows that scaling of $\mu$ is superlinear with $\norm{A}_1$ for $\tau=10^{-8}$. We thus settle for a numerical determination of the scaling relation between $\mu$ and $d$ for $\tau=10^{-8}$ by increasing the dimensional cutoff for both fluxonium qubits simultaneously. The observed scaling is $\mu\sim \Theta (d^{0.65})$, see Fig.\ \ref{fig:muscaling}(h). 


